\documentclass[aspectratio=169]{beamer}

% Theme selection
\usetheme{Madrid}
\usecolortheme{whale}

% Packages
\usepackage[utf8]{inputenc}
\usepackage{graphicx}
\usepackage{booktabs}
\usepackage{xcolor}
\usepackage{amssymb}

% Define custom colors
\definecolor{orange}{RGB}{255,140,0}
\definecolor{blue_dead}{RGB}{0,0,139} % Dark blue for "Dead"
\definecolor{green_triage}{RGB}{0,128,0}

% Meta Information
\title[Hospital Chatbot]{Hospital Chatbot for Patient Color-Coded Clinical Pathways}
\subtitle{AI-Driven Triage and Navigation for Resource-Constrained Settings}
\author{Project Synopsis}
\institute[]{Department of Mathematics}
\date{\today}

\begin{document}
	
	% Slide 1: Title
	\begin{frame}
		\titlepage
	\end{frame}
	
	% Slide 2: Introduction
	\begin{frame}{1. Introduction}
		\begin{block}{Context: Color-Coded Triage}
			Hospitals utilize standardized \textbf{color-coded triage systems} to classify patients by clinical severity. This ensures that critical resources are allocated to those who need them most.
		\end{block}
		
		\vspace{0.2cm}
		
		\begin{columns}[t]
			\begin{column}{0.48\textwidth}
				\textbf{The Triage Spectrum:}
				\begin{itemize}
					\item \textcolor{red}{\textbf{Red:}} Emergency (Immediate resuscitation).
					\item \textcolor{orange}{\textbf{Orange:}} Very Urgent (Potentially life-threatening).
					\item \textcolor{yellow}{\textbf{Yellow:}} Urgent (Serious but stable).
					\item \textcolor{green_triage}{\textbf{Green:}} Non-Urgent (Minor/Routine).
					\item \textcolor{blue_dead}{\textbf{Blue:}} Dead (Certified death on arrival).
				\end{itemize}
			\end{column}
			\begin{column}{0.48\textwidth}
				\begin{exampleblock}{The Innovation}
					An \textbf{AI-powered chatbot} acts as a "Virtual Triage Nurse," available 24/7. It assesses symptoms to assign these colors and directs patients to the correct department (e.g., Mortuary for Blue, Resuscitation for Red), streamlining flow before they even reach the desk.
				\end{exampleblock}
			\end{column}
		\end{columns}
	\end{frame}
	
	% Slide 3: Problem Statement
	\begin{frame}{2. Problem Statement}
		\begin{columns}[t]
			\begin{column}{0.48\textwidth}
				\begin{block}{The LMIC Context}
					In many \textbf{low- and middle-income countries}, healthcare systems face:
					\begin{itemize}
						\item \textbf{Severe Staff Shortages:} Low physician-to-patient ratios lead to bottlenecks.
						\item \textbf{Overcrowded Waiting Rooms:} Patients waiting hours without initial assessment.
						\item \textbf{Resource Misallocation:} Highly skilled staff spending time on minor cases while emergencies wait.
					\end{itemize}
				\end{block}
			\end{column}
			
			\begin{column}{0.48\textwidth}
				\begin{alertblock}{Critical Consequences}
					\begin{enumerate}
						\item \textbf{Undetected Deterioration:} "Yellow" patients turning "Orange" or "Red" while waiting unnoticed in the general queue.
						\item \textbf{System Congestion:} Non-urgent ("Green") cases clogging the Emergency Department instead of visiting Outpatient Clinics.
						\item \textbf{Patient Confusion:} Lack of clarity on where to go leads to anxiety and patients leaving without care (LWBS).
					\end{enumerate}
				\end{alertblock}
			\end{column}
		\end{columns}
		
		
	
	\end{frame}
	
	\begin{frame}{2.Problem Statement - Continued}
			\vspace{0.2cm}
		\textbf{Core Gap:} Lack of an automated, accessible system to identify a patient's \textbf{color code} and guide them immediately upon arrival.
	\end{frame}
	
	% Slide 4: Objectives
	\begin{frame}{3. Objectives}
		\begin{block}{General Objective}
			To develop an \textbf{AI-powered hospital chatbot} that assigns standardized triage colors (Red, Orange, Yellow, Green, Blue) and navigates patients through specific clinical pathways to reduce delays.
		\end{block}
		
		\vspace{0.09cm}
		\textbf{Specific Objectives:}
		\begin{itemize}
			\item \textbf{Algorithm Development:} Create logic based on validated protocols (e.g., SATS/TEWS) to calculate urgency colors accurately.
			\item \textbf{Pathway Mapping:} Define distinct workflows for each color:
			\begin{itemize} \footnotesize
				\item \emph{Red/Orange:} Immediate alert to ER staff.
				\item \emph{Yellow:} Direction to Urgent Care Zone.
				\item \emph{Green:} Direction to Pharmacy/OPD.
				\item \emph{Blue:} Protocol for sensitivity and transfer to Mortuary.
			\end{itemize}
			\item \textbf{Accuracy Evaluation:} Compare chatbot color assignments against human nurse triage to assess \textbf{sensitivity} and safety.
			
		\end{itemize}
	\end{frame}
	
	\begin{frame}{3. Objectives - Continued}
		\begin{itemize}
			\item \textbf{Impact Measurement:} Assess reduction in \textbf{Time to Provider} and overall waiting room congestion.
		\end{itemize}
		
	\end{frame}
	
	% Slide 5: Methodology
	\begin{frame}{4. Methodology}
		\begin{columns}[t]
			\begin{column}{0.48\textwidth}
				\begin{block}{Study Design \& Setting}
					\begin{itemize}
						\item \textbf{Design:} Mixed-methods approach.
						\begin{itemize} \footnotesize
							\item Phase 1: Agile Development of the Chatbot.
							\item Phase 2: Comparative Accuracy Assessment.
						\end{itemize}
						\item \textbf{Setting:} District or Regional Hospital in an LMIC setting utilizing standard color-coded triage.
					\end{itemize}
				\end{block}
			\end{column}
			
			\begin{column}{0.48\textwidth}
				\begin{block}{Data Collection \& Analysis}
					\begin{itemize}
						\item \textbf{Algorithm Validation:} Retrospective testing using historical patient vitals to check if the AI predicts the same color as the doctors did.
						\item \textbf{Usability Testing:} Surveys (SUS Scale) with patients to ensure the chatbot is easy to use for non-tech-savvy users.
						\item \textbf{Key Metrics:}
						\begin{itemize} \footnotesize
							\item Sensitivity (catching sick patients).
							\item Reduction in Waiting Times.
							\item Rate of "Left Without Being Seen".
						\end{itemize}
					\end{itemize}
				\end{block}
			\end{column}
		\end{columns}
	\end{frame}
	
	% Slide 6: Significance
	\begin{frame}{5. Significance}
		Implementation of the ``Virtual Triage Nurse'' aims to transform hospital flow:
		
		\vspace{0.3cm}
		
		\begin{enumerate}
			\item \textbf{Decongest the Emergency Department:} \\
			By accurately routing \textcolor{green_triage}{\textbf{Green}} patients to clinics and \textcolor{blue_dead}{\textbf{Blue}} cases to appropriate facilities, the ER beds remain open for true emergencies.
			
			\item \textbf{Improve Patient Safety:} \\
			Immediate flagging of \textcolor{red}{\textbf{Red}}/\textcolor{orange}{\textbf{Orange}} patients ensures they bypass the queue, reducing mortality associated with waiting delays.
			
			\item \textbf{Standardize Care Delivery:} \\
			Providing consistent, algorithmic assessment 24/7 removes human variability and fatigue from the initial sorting process.
		\end{enumerate}
	\end{frame}
	
\end{document}